% =====================================================================
% Predicción de la duración de desvíos de ruta · Lab. de IA · TP1
% Presentación del Grupo 5 (Sebastián, Bautista y Davor) — Tecfield
%
% Compilar con pdfLaTeX (en Overleaf: menú > Compiler > pdfLaTeX).
% Archivos del proyecto: presentacion.tex (este), tema-tp.sty (diseño),
% guion.tex (texto del orador), logo-uncoma.png, logo-informatica.png.
%
% Guion del orador (texto de cada lámina, tomado del video): descomentar la
% línea de abajo para generar un PDF solo con las notas, y volver a comentarla
% para el PDF de la presentación:
%   \setbeameroption{show only notes}
% =====================================================================
\documentclass[aspectratio=169,10pt]{beamer}
\usepackage{tema-tp}
% Guion del orador. Uso: \g{escena}{beat} dentro de \note{...}. Editable desde Overleaf.
\makeatletter
\newcommand{\g}[2]{\csname guion@#1@#2\endcsname\ }

\expandafter\def\csname guion@intro@hook\endcsname{Un vehículo de Tecfield se sale de su ruta. Llega el primer aviso, y la pregunta es una sola: ¿cuánto va a durar este desvío?}
\expandafter\def\csname guion@intro@tp\endcsname{Ese es nuestro trabajo práctico de Laboratorio de Inteligencia Artificial. Grupo 5: Sebastián, Bautista y Davor.}
\expandafter\def\csname guion@intro@mapa\endcsname{Vamos a recorrer el problema, los datos, el corpus, los modelos y los resultados.}

\expandafter\def\csname guion@problema@flota\endcsname{Tecfield opera flotas que llevan personal a obras y yacimientos. Cada vehículo tiene una hoja de ruta, con el recorrido planificado.}
\expandafter\def\csname guion@problema@detector\endcsname{Un detector compara la posición GPS contra ese recorrido. Cuando el vehículo se aleja más de 500 m, empieza a registrar lecturas de desvío. Hoy todas las alertas se atienden con la misma prioridad.}
\expandafter\def\csname guion@problema@t0\endcsname{Al primer punto fuera de ruta lo llamamos T\textsubscript{0}. Y en T\textsubscript{0} nadie sabe si el vehículo tomó una calle alternativa y vuelve en un minuto, o si va a quedar afuera una hora.}
\expandafter\def\csname guion@problema@etiquetas\endcsname{No hay etiquetas que digan por qué pasó cada desvío. Pero sí queda registrado cuánto duró: la última lectura menos la primera de cada evento, en segundos.}
\expandafter\def\csname guion@problema@objetivo\endcsname{Entonces el objetivo es una regresión: en T\textsubscript{0}, y solo con lo que se sabe hasta ese momento, predecir cuántos segundos va a durar el desvío. Una alerta temprana.}

\expandafter\def\csname guion@datos@tres\endcsname{Tecfield nos pasó tres colecciones exportadas de MongoDB: casi 2 GB de JSON.}
\expandafter\def\csname guion@datos@desvios\endcsname{Desvíos hoja de ruta: las lecturas del detector. De acá salen los eventos y la duración.}
\expandafter\def\csname guion@datos@hojas\endcsname{Hoja de ruta: el plan del viaje, con paradas, pasajeros y el trazado de OSRM, que da el largo y la duración del recorrido.}
\expandafter\def\csname guion@datos@rastreo\endcsname{Y rastreo satelital: el GPS crudo de cada vehículo, que usamos para ver qué venía haciendo antes de T\textsubscript{0}.}
\expandafter\def\csname guion@datos@trampas\endcsname{Al explorarlas aparecieron trampas. Las coordenadas vienen invertidas, las fechas vienen en hora universal, y el vehículo tiene un nombre distinto en cada colección.}
\expandafter\def\csname guion@datos@vigencia\endcsname{Además, el 61\,\% de los desvíos cae fuera de la vigencia de su hoja de ruta, y un mismo punto se guarda varias veces, contra hojas distintas.}

\expandafter\def\csname guion@corpus@intro\endcsname{Ahora, el corazón del trabajo: el corpus. Una tabla con una fila por desvío, los atributos que se conocen en T\textsubscript{0}, y la duración.}
\expandafter\def\csname guion@corpus@p1\endcsname{Se arma en 7 pasos. Primero, se tiran las lecturas inválidas: posición [0, 0], o fechas fuera de 2025 y 2026.}
\expandafter\def\csname guion@corpus@p2\endcsname{Segundo, se agrupa por evento, y T\textsubscript{0} es el primer ping. De ese ping salen la distancia hasta la ruta, la velocidad y la posición.}
\expandafter\def\csname guion@corpus@p3\endcsname{Tercero, se cruza con la hoja de ruta. Cuarto, se sacan los repetidos: uno por vehículo y T\textsubscript{0}, prefiriendo la hoja vigente.}
\expandafter\def\csname guion@corpus@p5\endcsname{Quinto, el rastreo previo: el GPS de los 10 minutos antes de T\textsubscript{0}, sin tocar nada del propio T\textsubscript{0}.}
\expandafter\def\csname guion@corpus@p6\endcsname{Sexto, la hora, el día y la franja. Y séptimo, el filtro final: dos lecturas o más, duración mayor a cero, y menos de 4 horas.}

\expandafter\def\csname guion@embudo@inicio\endcsname{Arrancamos con 37.720 eventos.}
\expandafter\def\csname guion@embudo@repetidos\endcsname{Se van 2.007 repetidos.}
\expandafter\def\csname guion@embudo@cortos\endcsname{19.240, con una sola lectura o duración cero.}
\expandafter\def\csname guion@embudo@tope\endcsname{Y 2.841 que llegan al tope de 4 horas: no se tiran, quedan aparte.}
\expandafter\def\csname guion@embudo@quedan\endcsname{Quedan 13.632 eventos, de 556 vehículos. Con eso se entrena.}

\expandafter\def\csname guion@atributos@grupos\endcsname{Los atributos vienen de tres lugares.}
\expandafter\def\csname guion@atributos@ping\endcsname{Del primer ping: distancia a la ruta, velocidad, posición, hora, día y franja.}
\expandafter\def\csname guion@atributos@hoja\endcsname{De la hoja de ruta: tipo de vehículo y de hoja, obra, largo y duración del recorrido, plazas, pasajeros, y si la hoja estaba vigente.}
\expandafter\def\csname guion@atributos@previo\endcsname{Y del rastreo previo: cantidad de puntos, velocidad, desplazamiento y rumbo de los 10 minutos anteriores.}

\expandafter\def\csname guion@evaluacion@baseline\endcsname{Para saber si un modelo sirve, hace falta una referencia. La nuestra es la mediana: predecir siempre la misma duración, sin mirar nada. Así se equivoca, en promedio, 32,1 minutos.}
\expandafter\def\csname guion@evaluacion@metricas\endcsname{Medimos tres métricas, en minutos. El MAE es el error absoluto medio. El error mediano es el del evento del medio. Y el RMSE, que se agregó después de la reunión con Tecfield, pesa más los errores grandes: en el ejemplo da 30 contra 22 de MAE.}
\expandafter\def\csname guion@evaluacion@grupos\endcsname{La validación cruzada es de 5 particiones, por vehículo. Un mismo vehículo nunca está en entrenamiento y en prueba a la vez, así que el modelo no memoriza un móvil que ya vio.}
\expandafter\def\csname guion@evaluacion@tres\endcsname{Contra esa referencia comparamos tres modelos: un árbol de decisión, un Random Forest, y XGBoost, que es el que nos recomendó Tecfield.}

\expandafter\def\csname guion@arbol@real\endcsname{El primer modelo es el árbol de decisión: hace preguntas sobre los atributos, y cada respuesta lo lleva por una rama hasta una hoja, que predice el promedio de sus eventos. En nuestro árbol, la primera pregunta es si el desvío trae el tipo de hoja informado. Si lo trae, la mediana es de 40 minutos. Si no, de 14. Y después, en las dos ramas, pregunta por la distancia a la ruta. Lo limitamos a 8 niveles para que no memorice.}
\expandafter\def\csname guion@arbol@resultado\endcsname{Resultado: 28,8 minutos de MAE y 47,5 de RMSE, un 10\,\% mejor que la mediana. La ventaja es que se puede dibujar y explicar. El límite, que un solo árbol es inestable: un cambio chico en los datos le cambia la estructura.}

\expandafter\def\csname guion@bosque@que\endcsname{El segundo modelo es el Random Forest: en vez de un árbol, 200 árboles.}
\expandafter\def\csname guion@bosque@azar\endcsname{Cada árbol se entrena con una muestra distinta de eventos, y en cada pregunta mira un subconjunto distinto de atributos. Así, cada uno se equivoca de una manera diferente.}
\expandafter\def\csname guion@bosque@promedio\endcsname{La predicción final es el promedio de los doscientos, y al promediar esos errores se cancelan en parte.}

\expandafter\def\csname guion@xgb@que\endcsname{El tercer modelo es XGBoost, el que nos recomendaron en Tecfield. No promedia árboles independientes: los entrena de a uno, y cada árbol nuevo corrige los errores de los anteriores. Usamos 300 árboles chicos.}
\expandafter\def\csname guion@xgb@config\endcsname{Lo probamos con una primera configuración, sin ajustar, igual que el Random Forest.}

\expandafter\def\csname guion@importancia@tres\endcsname{Sobre las variables importantes: en los tres modelos la número uno es el tipo de hoja, con entre el 29 y el 51\,\% de la importancia. Después viene la distancia hasta la ruta, la posición y la velocidad.}
\expandafter\def\csname guion@importancia@rastreo\endcsname{El rastreo previo pesa poco, entre 3 y 10\,\%. Ojo: esto dice qué usa el modelo, no qué causa la duración.}

\expandafter\def\csname guion@resultados@tabla\endcsname{Lado a lado: la mediana se equivoca 32,1 minutos de MAE, el árbol 28,8, el Random Forest 27 y XGBoost 27,3. El RMSE muestra lo mismo, y es bastante más alto porque los desvíos muy largos pesan más. Random Forest y XGBoost casi empatan.}
\expandafter\def\csname guion@resultados@grupos\endcsname{Si sumamos los grupos de atributos de a uno, con el Random Forest: solo con el primer ping, 29,4. Con la hoja de ruta baja a 27. Y el rastreo previo ya no mueve nada.}
\expandafter\def\csname guion@resultados@porque\endcsname{¿Por qué? Porque falta GPS previo en el 42\,\% de los eventos.}

\expandafter\def\csname guion@cierre@conclusion\endcsname{Conclusión: los modelos le ganan a la mediana, pero por poco. Hay señal, sobre todo en el tipo de hoja y el primer ping, aunque el error sigue siendo grande.}
\expandafter\def\csname guion@cierre@repo\endcsname{El código, el notebook y los informes están en el repositorio. Gracias.}

\makeatother

%\setbeameroption{show only notes}

% ------------------------------------------------------- Ayudas de dibujo
% Barra rectangular: \barra{x}{y}{ancho}{alto}{color}
\newcommand{\tmph}{h}
\newcommand{\avatar}[1]{\tikz[baseline=(c.base)]\node[circle,fill=tpblue,text=white,inner sep=0pt,minimum size=0.5cm,font=\bfseries\tpTiny](c){#1};}
\newcommand{\barra}[5]{\path[fill=#5,rounded corners=2.5pt] (#1,#2) rectangle (#1+#3,#2+#4);}
% Nodo de árbol: \tn[opciones]{cx}{cy}{ancho}{alto}{texto}
\newcommand{\tn}[6][]{\node[fill=tpcard,draw=tpborderhi,rounded corners=4pt,line width=0.6pt,
  minimum width=#4cm,minimum height=#5cm,align=center,inner sep=2pt,text=tptext,#1] at (#2,#3) {#6};}
% Arbolito de adorno: \minitree{x}{y}{escala}{color}
\newcommand{\minitree}[4]{\begin{scope}[shift={(#1,#2)},scale=#3]
  \draw[#4!60,line width=0.6pt] (0,0)--(-0.4,0.4) (0,0)--(0.4,0.4) (-0.4,0.4)--(-0.6,0.8) (-0.4,0.4)--(-0.2,0.8) (0.4,0.4)--(0.2,0.8) (0.4,0.4)--(0.6,0.8);
  \fill[#4] (0,0) circle (0.07) (-0.4,0.4) circle (0.06) (0.4,0.4) circle (0.06);
  \fill[#4!55] (-0.6,0.8) circle (0.05) (-0.2,0.8) circle (0.05) (0.2,0.8) circle (0.05) (0.6,0.8) circle (0.05);
\end{scope}}

% Ruta con desvío (diagrama del problema). \rutaplot{x}{y}{escala}
\newcommand{\rutapath}{(0,2.7) .. controls (1.2,2.8) and (2.0,1.7) .. (3.6,1.6) .. controls (5.2,1.5) and (5.9,0.6) .. (7.8,0.4)}
\newcommand{\rutaplot}[3]{\begin{scope}[shift={(#1,#2)},scale=#3]
  \draw[line width=18pt,tpskyhi!85] \rutapath;
  \draw[line width=2.2pt,tpblue] \rutapath;
  \draw[tpbad,line width=1.5pt] plot[smooth] coordinates {(4.0,1.58) (4.4,2.25) (5.0,2.7) (5.7,2.75) (6.3,2.2) (6.6,1.3) (6.75,0.62)};
  \foreach \p in {(4.4,2.25),(5.0,2.7),(5.7,2.75),(6.3,2.2),(6.6,1.3)}{\fill[tpbad] \p circle (0.075);}
  \draw[tpbad,line width=0.9pt] (4.4,2.25) circle (0.17);
  \draw[tpink,line width=1.1pt,fill=tpcard] (0,2.7) circle (0.13); \fill[tpink] (0,2.7) circle (0.05);
  \fill[tpbluedeep,rounded corners=1.5pt] (7.68,0.28) rectangle (7.92,0.52);
  \node[anchor=south,text=tpbluedeep] at (7.8,0.2) {\lab{Obra}};
  \node[anchor=north west,text=tpmuted] at (-0.1,2.95) {\lab{Vehículo}};
  \node[anchor=north west,text=tpbluedeep] at (1.7,3.0) {\lab{Hoja de ruta}};
  \node[anchor=south west,text=tpmuted] at (1.0,1.45) {\lab{Banda de 500 m}};
  \node[anchor=west,fill=tpbadbg,draw=tpbad!60,rounded corners=2pt,inner xsep=3pt,inner ysep=1.5pt,text=tpbad,font=\tpTiny\bfseries] at (4.75,3.1) {\T\ · primer ping};
  \draw[tpbad!60,line width=0.5pt] (4.4,2.42) -- (4.75,3.0);
\end{scope}}

\begin{document}

% =====================================================================
% 1 · PORTADA
% =====================================================================
\begin{frame}[plain]
\begin{lamina}
  % ruta de fondo
  \draw[line width=15pt,tpskyhi!70] (-0.3,8.0) .. controls (2,8.2) and (3,7.0) .. (5.5,7.35) .. controls (8,7.7) and (9,6.7) .. (11,6.9) .. controls (13,7.1) and (14.5,6.6) .. (16.3,6.4);
  \draw[line width=2pt,tpblue] (-0.3,8.0) .. controls (2,8.2) and (3,7.0) .. (5.5,7.35) .. controls (8,7.7) and (9,6.7) .. (11,6.9) .. controls (13,7.1) and (14.5,6.6) .. (16.3,6.4);
  \draw[tpbad,line width=1.4pt] plot[smooth] coordinates {(5.5,7.35) (6.2,8.0) (7.3,8.35) (8.3,8.1) (8.9,7.4) (9.2,6.95)};
  \foreach \p in {(6.2,8.0),(7.3,8.35),(8.3,8.1),(8.9,7.4)}{\fill[tpbad] \p circle (0.07);}
  \draw[tpink,fill=tpcard,line width=1.1pt] (0.9,8.05) circle (0.13); \fill[tpink] (0.9,8.05) circle (0.05);
  \fill[tpbluedeep,rounded corners=1.5pt] (15.0,6.42) rectangle (15.25,6.67);

  \node[anchor=north west] at (0.9,0.55) {\includegraphics[height=1.1cm]{logo-informatica.png}};
  \node[anchor=north east] at (15.1,0.55) {\includegraphics[height=1.1cm]{logo-uncoma.png}};
  \node[align=center,text=tpblue] at (8,2.1) {\lab{Facultad de Informática · UNComa}};
  \node[align=center,text=tpink,text width=14cm] at (8,3.35) {\bfseries\tpHero Laboratorio de\\Inteligencia Artificial};
  \node[align=center,text=tpmuted,text width=14cm] at (8,4.75) {\tpHthree TP1 · Predicción de la duración de desvíos de ruta};
  \node at (8,5.5) {\chipsolid{Grupo 5}\hspace{4pt}\chip{Sebastián}\hspace{4pt}\chip{Bautista}\hspace{4pt}\chip{Davor}};
  \node[text=tpdim,font=\ttfamily\tpTiny] at (8,6.05) {\textls[60]{CON TECFIELD · CURSADO 2026}};
\end{lamina}
\note{\g{intro}{hook} \g{intro}{tp}}
\end{frame}

% =====================================================================
% 2 · QUÉ VAMOS A RECORRER
% =====================================================================
\begin{frame}
\tpsec{Qué vamos a recorrer}{}
\begin{lamina}
  \eyebrow{0.9}{1.4}{Mapa de la presentación}
  \pbig{0.9}{1.75}{Cinco partes, \hl{tres voces}}
  \draw[tpborderhi,line width=1.2pt] (2.2,3.9) -- (13.8,3.9);
  \foreach \i/\t in {1/El problema,2/Los datos,3/El corpus,4/Los modelos,5/Resultados}{%
    \pgfmathsetmacro{\x}{2.2+(\i-1)*2.9}
    \num{\x}{3.9}{\i}
    \node[font=\bfseries\tpSmall,text=tpink] at (\x,4.5) {\t};}
  \card{0.9}{5.1}{4.5}{2.9}{%
    \avatar{B}\ {\bfseries\tpSmall Bautista}\par\vspace{5pt}
    \lab{Parte 1}\par\vspace{3pt}{\bfseries\tpHthree\hl{El corpus}}\par\vspace{3pt}{\tpSmall\color{tpmuted}De 37.720 eventos a una tabla para entrenar}}
  \card{5.75}{5.1}{4.5}{2.9}{%
    \avatar{S}\ {\bfseries\tpSmall Sebastián}\par\vspace{5pt}
    \lab{Parte 2}\par\vspace{3pt}{\bfseries\tpHthree\hl{Árbol de decisión}}\par\vspace{3pt}{\tpSmall\color{tpmuted}Modelo 1 · interpretable}}
  \card{10.6}{5.1}{4.5}{2.9}{%
    \avatar{D}\ {\bfseries\tpSmall Davor}\par\vspace{5pt}
    \lab{Parte 3}\par\vspace{3pt}{\bfseries\tpHthree\hl{Random Forest y XGBoost}}\par\vspace{3pt}{\tpSmall\color{tpmuted}Modelos 2 y 3 · bosque y boosting}}
\end{lamina}
\note{\g{intro}{mapa} \g{intro}{reparto}}
\end{frame}

% =====================================================================
% 3 · EL PROBLEMA
% =====================================================================
\begin{frame}
\tpsec{El problema}{}
\begin{lamina}
  \eyebrow{0.9}{1.3}{Tecfield · transporte de personal}
  \card{0.9}{1.65}{5.2}{1.4}{\lab{La flota}\par\vspace{2pt}\tpSmall Tecfield lleva personal a obras y yacimientos}
  \card{0.9}{3.15}{5.2}{1.4}{\lab{Hoja de ruta}\par\vspace{2pt}\tpSmall Cada vehículo tiene su recorrido planificado}
  \card{0.9}{4.65}{5.2}{1.4}{\lab{Detector}\par\vspace{2pt}\tpSmall Compara la posición GPS contra ese recorrido}
  \cardhi{0.9}{6.15}{5.2}{1.4}{\lab{Desvío}\par\vspace{2pt}\tpSmall A más de 500 m de la ruta registra lecturas de desvío}
  \card{6.4}{1.65}{8.7}{5.9}{}
  \rutaplot{6.95}{3.1}{1.0}
  \cardhi{0.9}{7.7}{14.2}{1.0}{\tpSmall Al primer punto fuera de ruta lo llamamos \hl{\T}. En ese instante, nadie sabe cuánto va a durar el desvío.}
\end{lamina}
\note{\g{problema}{flota} \g{problema}{detector} \g{problema}{t0}}
\end{frame}

% =====================================================================
% 5 · OBJETIVO: UNA REGRESIÓN
% =====================================================================
\begin{frame}
\tpsec{El problema}{}
\begin{lamina}
  \eyebrow{0.9}{1.3}{Objetivo}
  \pbig{0.9}{1.6}{Una \hl{regresión}}
  \card{0.9}{2.65}{6.9}{2.5}{%
    \lab{No hay etiquetas del motivo}\par\vspace{5pt}
    {\tpSmall\color{tpdim}%
    \textcircled{\tiny1} \sout{Falla de GPS}\par\vspace{2pt}
    \textcircled{\tiny2} \sout{Hoja de ruta mal configurada}\par\vspace{2pt}
    \textcircled{\tiny3} \sout{Vehículo mal asignado}\par\vspace{2pt}
    \textcircled{\tiny4} \sout{Recorrido distinto}}}
  \cardc{tpok}{8.2}{2.65}{6.9}{2.5}{%
    \lab{Lo que sí queda registrado}\par\vspace{5pt}
    {\bfseries\tpHthree Cuánto duró el desvío}\par\vspace{2pt}
    {\tpSmall\color{tpmuted}Última lectura $-$ primera, en segundos.}}
  \draw[tpok,line width=1.4pt] (8.8,4.35) -- (14.5,4.35);
  \fill[tpbad] (8.8,4.35) circle (0.09); \fill[tpbad] (14.5,4.35) circle (0.09);
  \node[anchor=north west,text=tpmuted,font=\ttfamily\tpTiny] at (8.6,4.55) {PRIMER PING};
  \node[anchor=north east,text=tpmuted,font=\ttfamily\tpTiny] at (14.7,4.55) {ÚLTIMO PING};
  \cardhi{0.9}{5.4}{6.4}{1.6}{\lab{Entrada · en \T}\par\vspace{2pt}{\tpSmall\color{tpmuted}Solo lo que se sabe hasta ese momento}\par\vspace{4pt}\chip{Primer ping}\ \chip{Hoja de ruta}\ \chip{Rastreo previo}}
  \draw[flecha] (7.4,6.2) -- (8.2,6.2);
  \node[fill=tpnavy,rounded corners=6pt,minimum width=1.5cm,minimum height=1.3cm,text=white,align=center,font=\bfseries\tpSmall] at (9.0,6.2) {$f(x)$\\[-1pt]{\mdseries\tpTiny Modelo}};
  \draw[flecha] (9.8,6.2) -- (10.6,6.2);
  \cardhi{10.7}{5.4}{4.4}{1.6}{\lab{Salida}\par\vspace{2pt}{\bfseries\tpBody Duración del desvío}\par\vspace{3pt}{\bfseries\tpHthree\hl{¿?} segundos}}
  \path[fill=tpblue,rounded corners=6pt] (0.9,7.3) rectangle (15.1,8.1);
  \node[anchor=west,text=white,font=\bfseries\tpBody] at (1.2,7.7) {Alerta temprana};
  \node[anchor=east,text=white!85,font=\tpSmall] at (14.8,7.7) {En el primer ping, estimar cuánto va a durar.};
\end{lamina}
\note{\g{problema}{etiquetas} \g{problema}{objetivo}}
\end{frame}

% =====================================================================
% 7 · LOS DATOS
% =====================================================================
\begin{frame}
\tpsec{Los datos}{}
\begin{lamina}
  \eyebrow{0.9}{1.3}{Tres colecciones · exportadas de MongoDB · casi 2 GB de JSON}
  \pbig{0.9}{1.6}{Los \hl{datos} de Tecfield}
  \card{0.9}{2.65}{4.55}{2.6}{\tmono{desviosHojaRuta}\par\vspace{2pt}\chip{98 MB}\par\vspace{4pt}\tpSmall Las lecturas del detector. De acá salen los \textbf{eventos} y la \textbf{duración}.}
  \card{5.75}{2.65}{4.55}{2.6}{\tmono{hojaDeRuta}\par\vspace{2pt}\chip{1,5 GB}\par\vspace{4pt}\tpSmall El plan del viaje: paradas, pasajeros y el trazado OSRM del recorrido.}
  \card{10.6}{2.65}{4.5}{2.6}{\tmono{rastreoSatelital}\par\vspace{2pt}\chip{284 MB}\par\vspace{4pt}\tpSmall El GPS crudo de cada vehículo: lo que venía haciendo \textbf{antes} de \T.}
  \node[txt,text=tpwarn] at (0.9,5.5) {\lab{Trampas que aparecieron al explorar}};
  \foreach \i/\a/\b in {0/Coordenadas invertidas/{Desvíos: \texttt{[lat, lon]}. GPS: \texttt{[lon, lat]}},
                        1/Fechas en hora universal/{Se pasan a la hora de Argentina},
                        2/Tres nombres para el vehículo/{\texttt{flotaId}, \texttt{flotaID}, \texttt{vehiculo.id}},
                        3/Hojas vencidas/{El 61\,\% de los desvíos cae fuera de la vigencia de su hoja},
                        4/Desvíos repetidos/{Un mismo punto se guarda contra varias hojas}}{%
    \pgfmathsetmacro{\x}{0.9+\i*2.875}
    \cardc{tpwarn}{\x}{5.85}{2.7}{2.35}{{\bfseries\tpSmall\a}\par\vspace{3pt}{\tpTiny\color{tpmuted}\b}}}
\end{lamina}
\note{\g{datos}{tres} \g{datos}{desvios} \g{datos}{hojas} \g{datos}{rastreo} \g{datos}{trampas} \g{datos}{vigencia}}
\end{frame}

% =====================================================================
% 9 · CÓMO SE ARMA EL CORPUS (Bautista)
% =====================================================================
\begin{frame}
\tpsec{Cómo se arma el corpus}{}
\begin{lamina}
  \eyebrow{0.9}{1.3}{Una fila por desvío · atributos en \T\ y la duración}
  \pbig{0.9}{1.6}{Siete pasos para armar el \hl{corpus}}
  \foreach \i/\t/\d in {1/Lecturas válidas/{Se tiran las de posición \texttt{[0, 0]} o con fecha fuera de 2025--2026.},
                        2/Un registro por evento/{Se agrupa por \tmono{grupoDesvioId}. \T\ es el primer ping: distancia, velocidad y posición.},
                        3/Hoja de ruta/{Se cruza con la hoja: obra, tipo de vehículo, largo y duración OSRM, vigencia.},
                        4/Sin repetidos/{Uno por vehículo y \T, prefiriendo la hoja vigente.}}{%
    \pgfmathsetmacro{\y}{2.65+(\i-1)*1.4}
    \num{1.2}{\y+0.3}{\i}
    \node[txt,text width=6.4cm] at (1.7,\y) {{\bfseries\tpBody\t}\par\vspace{2pt}{\tpSmall\color{tpmuted}\d}};}
  \foreach \i/\t/\d in {5/Rastreo previo/{El GPS de los 10 minutos antes de \T\ (\tmono{join_asof}), sin tocar el propio \T.},
                        6/Fecha y hora/{Hora, día de la semana y franja del día.},
                        7/Filtro final/{2 lecturas o más, duración mayor a 0 y menos de 4 horas.}}{%
    \pgfmathsetmacro{\y}{2.65+(\i-5)*1.4}
    \num{8.7}{\y+0.3}{\i}
    \node[txt,text width=5.9cm] at (9.2,\y) {{\bfseries\tpBody\t}\par\vspace{2pt}{\tpSmall\color{tpmuted}\d}};}
  \cardc{tpok}{8.5}{6.65}{6.6}{1.45}{\lab{Sin fugas}\par\vspace{2pt}\tpSmall Solo se usa información de \T\ o anterior. Nada del final del desvío.}
\end{lamina}
\note{\g{corpus}{intro} \g{corpus}{p1} \g{corpus}{p2} \g{corpus}{p3} \g{corpus}{p5} \g{corpus}{p6}}
\end{frame}

% =====================================================================
% 10 · DE EVENTOS A CORPUS (Bautista)
% =====================================================================
\begin{frame}
\tpsec{De eventos a corpus}{}
\begin{lamina}
  \eyebrow{0.9}{1.3}{Cuántos eventos quedan}
  \pbig{0.9}{1.6}{De 37.720 eventos a \hl{13.632}}
  % barras: N = 37.720 eventos ≙ 9,6 cm, desde x = 5,0
  \node[txt,text=tpink] at (0.9,2.7) {\bfseries\tpHthree 37.720};
  \node[txt,text=tpmuted] at (0.9,3.15) {\tpTiny eventos de partida};
  \barra{5.0}{2.75}{9.6}{0.5}{tpskyhi}
  \node[txt,text=tpmuted] at (0.9,3.85) {\bfseries\tpHthree$-$2.007};
  \node[txt,text=tpmuted] at (0.9,4.3) {\tpTiny repetidos};
  \barra{14.09}{3.9}{0.51}{0.5}{tpdim}
  \node[txt,text=tpmuted] at (0.9,5.0) {\bfseries\tpHthree$-$19.240};
  \node[txt,text=tpmuted] at (0.9,5.45) {\tpTiny una sola lectura o duración 0};
  \barra{9.19}{5.05}{4.9}{0.5}{tpdim}
  \node[txt,text=tpwarn] at (0.9,6.15) {\bfseries\tpHthree$-$2.841};
  \node[txt,text=tpwarn] at (0.9,6.6) {\tpTiny tope de 4 horas};
  \barra{8.47}{6.2}{0.72}{0.5}{tpwarn}
  \draw[flecha,tpwarn] (9.4,6.45) -- (9.95,6.45);
  \cardc{tpwarn}{10.05}{5.85}{5.05}{1.3}{\chipsolid[tpwarn]{Aparte}\ {\bfseries\tpSmall No se tiran}\par\vspace{1pt}{\tpTiny\color{tpmuted}El detector cortó: no hay retorno medido.}}
  \node[txt,text=tpblue] at (0.9,7.25) {\bfseries\tpHone 13.632};
  \node[txt,text=tpblue] at (0.9,8.0) {\bfseries\tpTiny eventos en el corpus};
  \barra{5.0}{7.3}{3.47}{0.7}{tpblue}
  \node[anchor=west,text=tpink,font=\bfseries\tpHthree] at (8.75,7.65) {de \hl{556} vehículos};
\end{lamina}
\note{\g{embudo}{inicio} \g{embudo}{repetidos} \g{embudo}{cortos} \g{embudo}{tope} \g{embudo}{quedan}}
\end{frame}

% =====================================================================
% 11 · ATRIBUTOS EN T0 (Bautista)
% =====================================================================
\begin{frame}
\tpsec{Atributos en \T}{}
\begin{lamina}
  \eyebrow{0.9}{1.3}{Atributos candidatos · todos se conocen en \T}
  \pbig{0.9}{1.6}{Los atributos vienen de \hl{tres lugares}}
  \cardhi{0.9}{2.65}{4.55}{3.45}{{\color{tpblue}\bfseries\tpHthree 1}\ {\bfseries\tpBody Primer ping}\par\vspace{1pt}{\tpTiny\color{tpmuted}Lo que mide el primer GPS fuera de ruta}\par\vspace{5pt}
    \chip{distancia\_m}\ \chip{velocidad}\ \chip{velocidad\_vacia}\ \chip{ruta\_cercana}\ \chip{lat}\ \chip{lon}\ \chip{hora}\ \chip{dia\_semana}\ \chip{franja}}
  \card{5.75}{2.65}{4.55}{3.45}{{\color{tpblue}\bfseries\tpHthree 2}\ {\bfseries\tpBody Hoja de ruta}\par\vspace{1pt}{\tpTiny\color{tpmuted}Lo que dice la hoja del viaje}\par\vspace{5pt}
    \chip{tipoVehiculo}\ \chip{tipoHDR}\ \chip{obra\_id}\ \chip{km\_ida}\ \chip{min\_ida}\ \chip{plazas}\ \chip{pasajeros}\ \chip{hoja\_vigente}}
  \card{10.6}{2.65}{4.5}{3.45}{{\color{tpblue}\bfseries\tpHthree 3}\ {\bfseries\tpBody Rastreo previo}\par\vspace{1pt}{\tpTiny\color{tpmuted}Qué hizo en los 10 min anteriores}\par\vspace{5pt}
    \chip{gps\_previos}\ \chip{vel\_media\_previa}\ \chip{vel\_ultima}\ \chip{min\_desde\_ultimo\_gps}\ \chip{desplaz\_previo\_m}\ \chip{rumbo\_previo}}
  \cardc{tpbad}{0.9}{6.35}{14.2}{1.85}{}
  \node[txt,text=tpbad] at (1.2,6.6) {\bfseries\tpHthree Fugas: nunca entran};
  \node[anchor=north east,text=tpbad!80,font=\tpSmall] at (14.8,6.65) {Sería hacer trampa};
  \foreach \i/\a/\b/\sx in {0/retorno.*/Cómo terminó el desvío/1,1/lecturas/Cuántas lecturas tuvo/1,2/GPS posterior a \T/Cualquier dato posterior/0}{%
    \pgfmathsetmacro{\x}{1.2+\i*4.65}
    \path[fill=tpcard,draw=tpbad!40,rounded corners=4pt,line width=0.5pt] (\x,7.15) rectangle (\x+4.3,8.0);
    \ifnum\sx=1
      \node[anchor=west,font=\ttfamily\tpSmall,text=tpbad] at (\x+0.2,7.4) {\sout{\a}};
    \else
      \node[anchor=west,font=\ttfamily\tpSmall,text=tpbad] at (\x+0.2,7.4) {\a};
    \fi
    \node[anchor=west,font=\tpTiny,text=tpmuted] at (\x+0.2,7.78) {\b};}
\end{lamina}
\note{\g{atributos}{grupos} \g{atributos}{ping} \g{atributos}{hoja} \g{atributos}{previo} \g{atributos}{fugas}}
\end{frame}

% =====================================================================
% 13 · CÓMO SE EVALÚA
% =====================================================================
\begin{frame}
\tpsec{Cómo se evalúa}{}
\begin{lamina}
  \eyebrow{0.9}{1.3}{Referencia, métricas y validación}
  \pbig{0.9}{1.6}{La referencia: \hl{la mediana}}
  % 1 · referencia
  \card{0.9}{2.6}{4.55}{3.75}{\lab{Referencia}\par\vspace{2pt}{\bfseries\tpSmall Predecir siempre la mediana}}
  \draw[tpblue,line width=1.2pt] (1.2,4.85) -- (5.2,4.85);
  \foreach \p in {(1.5,5.2),(1.8,4.95),(2.1,5.3),(2.4,4.5),(2.7,5.1),(3.0,3.9),(3.3,5.25),(3.6,4.95),(3.9,4.3),(4.2,5.3),(4.5,4.9),(4.8,3.75),(5.0,5.1)}{\fill[tpblue!70] \p circle (0.055);}
  \node[txt,text=tpink] at (1.2,5.4) {\bfseries\tpHone 32,1};
  \tx[mut]{3.05}{5.5}{2.3}{\tpTiny min de error medio (MAE) de la referencia}
  % 2 · métricas
  \card{5.75}{2.6}{4.55}{3.75}{\lab{Métricas}\par\vspace{2pt}{\bfseries\tpSmall MAE, RMSE y error mediano}\par{\tpTiny\color{tpmuted}5 eventos de ejemplo, en minutos}}
  \foreach \i/\e in {0/14,1/61,2/4,3/22,4/9}{%
    \pgfmathsetmacro{\y}{4.1+\i*0.28}
    \pgfmathsetmacro{\w}{\e/61*3.4}
    \barra{6.1}{\y}{\w}{0.2}{tpskyhi}}
  \draw[tpblue,line width=0.9pt] (6.1+22/61*3.4,4.0) -- (6.1+22/61*3.4,5.5);
  \draw[tpink,line width=0.9pt,dashed] (6.1+14/61*3.4,4.0) -- (6.1+14/61*3.4,5.5);
  \node[anchor=north west,text=tpblue,font=\ttfamily\tpTiny] at (6.1+22/61*3.4+0.08,5.5) {MAE = 22 (media)};
  \node[anchor=north west,text=tpink,font=\ttfamily\tpTiny] at (6.1,5.75) {mediano = 14 (el del medio)};
  \node[anchor=north west,text=tpbad,font=\ttfamily\tpTiny] at (6.1,6.0) {RMSE = 30 (pesa los grandes)};
  % 3 · validación
  \card{10.6}{2.6}{4.5}{3.75}{\lab{Validación cruzada}\par\vspace{2pt}{\bfseries\tpSmall 5 particiones por vehículo}}
  \foreach \f [count=\k from 0] in {0,3,1,4,2,0,2,4,1,3,0,4,3,1,2,1,2,0,3,4,4,0,1,2}{%
    \pgfmathtruncatemacro{\r}{int(\k/6)}
    \pgfmathtruncatemacro{\c}{mod(\k,6)}
    \pgfmathsetmacro{\x}{10.95+\c*0.62}
    \pgfmathsetmacro{\y}{3.85+\r*0.42}
    \ifcase\f\def\col{tpblue}\or\def\col{tpnavy}\or\def\col{tpblue!55}\or\def\col{tpmuted}\or\def\col{tpok!70}\fi
    \path[fill=\col,rounded corners=2pt] (\x,\y) rectangle (\x+0.48,\y+0.32);}
  \tx[mut]{10.9}{5.65}{3.9}{\tpTiny Nunca en entrenamiento y prueba a la vez.}
  % comparación
  \node[txt,text=tpmuted] at (0.9,6.65) {\lab{Contra esa referencia comparamos tres modelos}};
  \card{0.9}{7.0}{4.55}{1.2}{{\bfseries\tpBody Árbol de decisión}\par\vspace{1pt}{\tpTiny\color{tpmuted}Modelo 1}}
  \card{5.75}{7.0}{4.55}{1.2}{{\bfseries\tpBody Random Forest}\par\vspace{1pt}{\tpTiny\color{tpmuted}Modelo 2}}
  \card{10.6}{7.0}{4.5}{1.2}{{\bfseries\tpBody XGBoost}\par\vspace{1pt}{\tpTiny\color{tpmuted}Modelo 3 · nuevo}}
\end{lamina}
\note{\g{evaluacion}{baseline} \g{evaluacion}{metricas} \g{evaluacion}{grupos} \g{evaluacion}{tres}}
\end{frame}

% =====================================================================
% 15 · EL ÁRBOL REAL (Sebastián)
% =====================================================================
\begin{frame}
\tpsec{Árbol de decisión}{}
\begin{lamina}
  \eyebrow{0.9}{1.3}{Nuestro árbol · entrenado con los 13.632 eventos}
  \pbig{0.9}{1.6}{¿Trae el \hl{tipo de hoja}?}
  \draw[tpblue,line width=1.8pt] (5.8,3.6) -- (3.3,4.3);
  \draw[tpborderhi,line width=1.2pt] (5.8,3.6) -- (8.3,4.3);
  \draw[tpborderhi,line width=1.2pt] (3.3,5.5) -- (2.05,6.05) (3.3,5.5) -- (4.55,6.05) (8.3,5.5) -- (7.05,6.05) (8.3,5.5) -- (9.55,6.05);
  \node[fill=tpbg,text=tpbluedeep,font=\ttfamily\tpTiny,inner sep=1.5pt] at (4.4,3.9) {sí};
  \node[fill=tpbg,text=tpmuted,font=\ttfamily\tpTiny,inner sep=1.5pt] at (7.2,3.9) {no};
  \tn[draw=tpblue,line width=1pt,fill=tpsky]{5.8}{3.15}{4.2}{0.9}{{\bfseries\tpSmall ¿\tmono{tipoHDR} informado?}\\[1pt]{\tpTiny\color{tpmuted}13.632 eventos · mediana 20 min}}
  \tn{3.3}{4.9}{4.1}{1.2}{{\ttfamily\tpTiny CON TIPO · 6.828 EVENTOS}\\[1pt]{\bfseries\tpBody\hl{mediana 40 min}}\\[1pt]{\tpTiny ¿distancia $\leq$ 2.938 m?}}
  \tn{8.3}{4.9}{4.1}{1.2}{{\ttfamily\tpTiny SIN TIPO · 6.804 EVENTOS}\\[1pt]{\bfseries\tpBody\hl{mediana 14 min}}\\[1pt]{\tpTiny ¿distancia $\leq$ 1.888 m?}}
  \tn{2.05}{6.55}{2.3}{0.9}{{\tpTiny ¿velocidad $\leq$ 33,5?}\\[1pt]{\bfseries\tpTiny\hl{mediana 33 min}}}
  \tn{4.55}{6.55}{2.3}{0.9}{{\tpTiny ¿latitud $\leq$ $-$38,75?}\\[1pt]{\bfseries\tpTiny\hl{mediana 50 min}}}
  \tn{7.05}{6.55}{2.3}{0.9}{{\tpTiny ¿latitud $\leq$ $-$39,16?}\\[1pt]{\bfseries\tpTiny\hl{mediana 10 min}}}
  \tn{9.55}{6.55}{2.3}{0.9}{{\tpTiny ¿dist. $\leq$ 13.958 m?}\\[1pt]{\bfseries\tpTiny\hl{mediana 16 min}}}
  \foreach \x in {2.05,4.55,7.05,9.55}{\draw[tpborderhi,line width=0.9pt,dash pattern=on 2pt off 2pt] (\x,7.0) -- (\x,7.55);
    \node[text=tpdim,font=\tpBody] at (\x,7.75) {$\vdots$};}
  \node[anchor=north west,text=tpmuted,font=\ttfamily\tpTiny] at (0.9,8.15) {Y SIGUE HASTA 8 NIVELES DE PROFUNDIDAD};
  \card{11.4}{2.65}{3.7}{5.5}{\lab{Resultado}\par\vspace{4pt}
    {\bfseries\tpHone\hl{28,8}}\par{\tpTiny\color{tpmuted}min de MAE · RMSE 47,5}\par\vspace{4pt}
    \chipsolid[tpok]{10 \% mejor}\par\vspace{8pt}
    {\bfseries\tpSmall\color{tpok}Ventaja}\par{\tpTiny\color{tpmuted}Se puede dibujar y explicar.}\par\vspace{6pt}
    {\bfseries\tpSmall\color{tpwarn}Límite}\par{\tpTiny\color{tpmuted}Un solo árbol es inestable: un cambio chico en los datos le cambia la estructura.}}
\end{lamina}
\note{\g{arbol}{real} \g{arbol}{resultado}}
\end{frame}

% =====================================================================
% 17 · RANDOM FOREST: QUÉ ES (Davor)
% =====================================================================
\begin{frame}
\tpsec{Random Forest}{}
\begin{lamina}
  \eyebrow{0.9}{1.3}{Modelo 2}
  \pbig{0.9}{1.6}{Random \hl{Forest}}
  \tx[mut]{0.9}{2.6}{5.2}{\tpSmall Un bosque de árboles de decisión, no uno solo. La predicción es el promedio de todos.}
  \node[txt,text=tpblue] at (0.9,3.4) {\bfseries\tpHero 200};
  \node[txt,text=tpmuted] at (3.35,4.1) {\tpBody árboles};
  \card{0.9}{4.95}{5.2}{1.55}{\lab{Cada árbol, una muestra distinta}\par\vspace{2pt}\tpSmall Se entrena con eventos elegidos al azar, con repetición.}
  \card{0.9}{6.65}{5.2}{1.55}{\lab{Cada pregunta, atributos distintos}\par\vspace{2pt}\tpSmall En cada pregunta mira solo un subconjunto de los atributos.}
  \card{6.4}{1.7}{8.7}{3.0}{}
  \foreach \r in {0,1,2}{\foreach \c in {0,...,7}{%
    \pgfmathsetmacro{\x}{7.1+\c*1.0}
    \pgfmathsetmacro{\y}{2.05+\r*0.85}
    \minitree{\x}{\y}{0.55}{tpblue}}}
  \node[anchor=south east,text=tpmuted,font=\ttfamily\tpTiny] at (14.9,4.6) {200 ÁRBOLES DISTINTOS};
  \card{6.4}{4.9}{8.7}{3.3}{}
  \node[txt,text=tpmuted] at (6.7,5.1) {\lab{Tabla de entrenamiento}};
  \node[anchor=north east,text=tpdim,font=\ttfamily\tpTiny] at (14.8,5.1) {FILA = EVENTO · COLUMNA = ATRIBUTO};
  \foreach \c/\n in {0/tipoHDR,1/distancia,2/lat,3/lon,4/velocidad,5/obra}{%
    \pgfmathsetmacro{\x}{7.4+\c*1.2}
    \node[font=\ttfamily\tpTiny,text=tpmuted] at (\x+0.5,5.65) {\n};}
  \foreach \r in {0,...,6}{\foreach \c in {0,...,5}{%
    \pgfmathsetmacro{\x}{7.4+\c*1.2}
    \pgfmathsetmacro{\y}{5.9+\r*0.3}
    \path[fill=tpsky,rounded corners=1.5pt] (\x,\y) rectangle (\x+1.0,\y+0.22);}}
  \foreach \r in {1,3,5}{\path[fill=tpblue,opacity=0.18,rounded corners=1.5pt] (7.3,5.9+\r*0.3-0.03) rectangle (14.7,5.9+\r*0.3+0.25);}
  \foreach \c in {1,2,4}{\path[draw=tpblue,line width=0.9pt,rounded corners=3pt] (7.35+\c*1.2,5.8) rectangle (8.45+\c*1.2,7.95);}
\end{lamina}
\note{\g{bosque}{que} \g{bosque}{azar} \g{bosque}{promedio}}
\end{frame}

% =====================================================================
% 18 · XGBOOST (Davor)
% =====================================================================
\begin{frame}
\tpsec{XGBoost}{}
\begin{lamina}
  \eyebrow{0.9}{1.3}{Modelo 3 · nuevo, recomendado por Tecfield}
  \pbig{0.9}{1.6}{\hl{XGBoost}: árboles que se corrigen}
  \tx[mut]{0.9}{2.6}{14}{\tpSmall A diferencia del bosque, los árboles no son independientes: se entrenan de a uno y cada uno corrige los errores del anterior.}
  \card{0.9}{3.5}{14.2}{2.6}{}
  \foreach \x/\l/\d in {2.3/Árbol 1/{Primera estimación},5.8/Árbol 2/{Corrige los errores del 1},9.3/Árbol 3/{Corrige lo que queda}}{%
    \minitree{\x}{4.0}{1.2}{tpblue}
    \node[text=tpink,font=\bfseries\tpSmall] at (\x,5.3) {\l};
    \node[text=tpmuted,font=\tpTiny] at (\x,5.65) {\d};}
  \foreach \x in {4.05,7.55}{\node[font=\bfseries\tpHone,text=tpdim] at (\x,4.5) {$+$};}
  \node[font=\bfseries\tpHone,text=tpdim] at (11.05,4.5) {$+\ \cdots$};
  \node[text=tpblue,font=\bfseries\tpHone] at (13.3,4.35) {300};
  \node[text=tpmuted,font=\tpTiny] at (13.3,4.85) {árboles chicos, en secuencia};
  \card{0.9}{6.35}{4.55}{1.85}{\lab{Bosque vs. boosting}\par\vspace{3pt}\tpSmall El bosque promedia árboles independientes. XGBoost los apila de a uno.}
  \card{5.75}{6.35}{4.55}{1.85}{\lab{Cómo lo configuramos}\par\vspace{3pt}{\ttfamily\tpTiny n\_estimators = 300\\learning\_rate = 0.05\\max\_depth = 4}}
  \cardc{tpwarn}{10.6}{6.35}{4.5}{1.85}{\lab{Ojo}\par\vspace{3pt}\tpSmall Es una primera configuración: no ajustamos los parámetros.}
\end{lamina}
\note{\g{xgb}{que} \g{xgb}{config}}
\end{frame}

% =====================================================================
% 19 · IMPORTANCIA DE VARIABLES
% =====================================================================
\begin{frame}
\tpsec{Importancia de variables}{}
\begin{lamina}
  \eyebrow{0.9}{1.3}{Qué usa cada modelo para predecir · la importancia de cada modelo suma 100 \%}
  \pbig{0.9}{1.6}{¿Qué variables pesan \hl{más}?}
  \card{0.9}{2.65}{9.6}{5.55}{}
  \foreach \n/\a/\b/\c [count=\i from 0] in {Tipo de hoja/0.5082/0.2920/0.4918,Distancia a la ruta/0.1399/0.1255/0.0473,Latitud/0.0678/0.0902/0.0303,Velocidad/0.0631/0.0776/0.0314,Hora/0.0440/0.0449/0.0328}{%
    \pgfmathsetmacro{\y}{3.0+\i*0.78}
    \pgfmathsetmacro{\wa}{\a/0.5082*3.4}
    \pgfmathsetmacro{\wb}{\b/0.5082*3.4}
    \pgfmathsetmacro{\wc}{\c/0.5082*3.4}
    \pgfmathsetmacro{\pa}{\a*100}
    \pgfmathsetmacro{\pb}{\b*100}
    \pgfmathsetmacro{\pc}{\c*100}
    \node[anchor=east,font=\tpSmall] at (4.3,\y+0.31) {\n};
    \barra{4.5}{\y}{\wa}{0.19}{tpdim}
    \barra{4.5}{\y+0.22}{\wb}{0.19}{tpblue}
    \barra{4.5}{\y+0.44}{\wc}{0.19}{tpnavy}
    \node[anchor=west,font=\ttfamily\tpTiny,text=tpmuted] at (4.5+\wa+0.1,\y+0.095) {\pgfmathprintnumber[fixed,precision=0]{\pa}\,\%};
    \node[anchor=west,font=\ttfamily\tpTiny,text=tpmuted] at (4.5+\wb+0.1,\y+0.315) {\pgfmathprintnumber[fixed,precision=0]{\pb}\,\%};
    \node[anchor=west,font=\ttfamily\tpTiny,text=tpmuted] at (4.5+\wc+0.1,\y+0.535) {\pgfmathprintnumber[fixed,precision=0]{\pc}\,\%};}
  \node[anchor=north west] at (1.2,7.6) {\tikz\fill[tpdim,rounded corners=1pt] (0,0) rectangle (0.25,0.18);\ {\tpTiny Árbol}\hspace{10pt}\tikz\fill[tpblue,rounded corners=1pt] (0,0) rectangle (0.25,0.18);\ {\tpTiny Random Forest}\hspace{10pt}\tikz\fill[tpnavy,rounded corners=1pt] (0,0) rectangle (0.25,0.18);\ {\tpTiny XGBoost}};
  \cardhi{10.8}{2.65}{4.3}{2.6}{\lab{En los tres modelos}\par\vspace{3pt}{\bfseries\tpBody El tipo de hoja es la variable número 1}\par\vspace{3pt}{\tpTiny\color{tpmuted}Entre el 29 y el 51 \% de la importancia. Después viene la distancia hasta la ruta.}}
  \card{10.8}{5.45}{4.3}{2.75}{\lab{Rastreo previo}\par\vspace{3pt}{\bfseries\tpBody Pesa poco: entre 3 y 10 \%}\par\vspace{3pt}{\tpTiny\color{tpmuted}Dice qué usa el modelo, no qué causa la duración.}}
\end{lamina}
\note{\g{importancia}{tres} \g{importancia}{rastreo}}
\end{frame}

% =====================================================================
% 20 · RESULTADOS LADO A LADO
% =====================================================================
\begin{frame}
\tpsec{Resultados}{}
\begin{lamina}
  \eyebrow{0.9}{1.3}{Validación cruzada por vehículo · 5 particiones · 13.632 eventos}
  \pbig{0.9}{1.6}{Lado a lado: \hl{el error de cada modelo}}
  \card{0.9}{2.65}{6.95}{4.6}{}
  \card{8.15}{2.65}{6.95}{4.6}{}
  \node[txt,text=tpmuted] at (1.2,2.9) {\lab{MAE · error medio · en minutos}};
  \node[txt,text=tpmuted] at (8.45,2.9) {\lab{RMSE · castiga los errores grandes}};
  \foreach \n/\m/\lm/\r/\lr/\c [count=\i from 0] in {Mediana/32.1/{32,1}/54.2/{54,2}/tpborderhi,Árbol de decisión/28.8/{28,8}/47.5/{47,5}/tpskyhi,Random Forest/27.0/{27,0}/45.8/{45,8}/tpblue,XGBoost/27.3/{27,3}/46.1/{46,1}/tpnavy}{%
    \pgfmathsetmacro{\y}{3.3+\i*0.92}
    \pgfmathsetmacro{\wm}{\m/32.1*4.2}
    \pgfmathsetmacro{\wr}{\r/54.2*4.2}
    \node[txt,font=\bfseries\tpSmall] at (1.2,\y) {\n};
    \node[txt,font=\bfseries\tpSmall] at (8.45,\y) {\n};
    \barra{1.2}{\y+0.28}{\wm}{0.42}{\c}
    \barra{8.45}{\y+0.28}{\wr}{0.42}{\c}
    \node[anchor=west,font=\ttfamily\small\bfseries,text=tpink] at (1.2+\wm+0.15,\y+0.49) {\lm};
    \node[anchor=west,font=\ttfamily\small\bfseries,text=tpink] at (8.45+\wr+0.15,\y+0.49) {\lr};}
  \cardhi{0.9}{7.5}{14.2}{0.95}{\tpSmall \textbf{Random Forest} baja el MAE un 16 \% respecto de la mediana y \textbf{XGBoost} un 15 \%: casi empatan (0,3 min de diferencia), con parámetros sin ajustar.}
\end{lamina}
\note{\g{resultados}{tabla}}
\end{frame}

% =====================================================================
% 22 · SUMANDO GRUPOS DE ATRIBUTOS
% =====================================================================
\begin{frame}
\tpsec{Resultados}{}
\begin{lamina}
  \eyebrow{0.9}{1.3}{Random Forest · MAE en minutos}
  \pbig{0.9}{1.6}{Sumando grupos de \hl{atributos}}
  \card{0.9}{2.65}{8.6}{5.55}{}
  % alto de barra = MAE / 32,1 * 3,3 ; base en y = 7,3
  \draw[tpdim,dashed,line width=0.7pt] (1.3,4.0) -- (9.1,4.0);
  \node[anchor=south east,text=tpmuted,font=\ttfamily\tpTiny] at (9.1,4.0) {MEDIANA · 32,1};
  \foreach \n/\val/\lbl/\p/\x/\c/\tc in {Primer ping/29.4/{29,4}/8 \%/1.6/tpskyhi/tpink,+ Hoja de ruta/27.0/{27,0}/16 \%/4.5/tpblue/white,+ Rastreo previo/27.0/{27,0}/16 \%/7.4/tpblue/white}{%
    \pgfmathsetmacro{\h}{\val/32.1*3.3}
    \barra{\x}{7.3-\h}{1.7}{\h}{\c}
    \node[anchor=north,font=\bfseries\tpHthree,text=\tc] at (\x+0.85,7.3-\h+0.15) {\lbl};
    \node[anchor=north,font=\ttfamily\tpTiny,text=\tc] at (\x+0.85,7.3-\h+0.6) {\p\ mejor};
    \node[anchor=north,font=\tpSmall] at (\x+0.85,7.4) {\n};}
  \card{9.8}{2.65}{5.3}{5.55}{\lab{¿Por qué el rastreo previo no ayuda?}}
  \draw[tpskyhi,line width=15pt] (12.45,4.9) circle (1.2);
  \draw[tpwarn,line width=15pt] (12.45,3.7) arc (-90:{-90+0.424*360}:1.2);
  \node[font=\bfseries\tpHone,text=tpink] at (12.45,4.8) {42 \%};
  \node[font=\tpTiny,text=tpmuted] at (12.45,5.25) {de los eventos};
  \tx{10.1}{6.5}{4.7}{\tpSmall\textcolor{tpwarn}{$\blacksquare$}\ \textbf{Sin GPS previo:} antes de \T\ no hay ninguna lectura del vehículo.}
\end{lamina}
\note{\g{resultados}{grupos} \g{resultados}{porque}}
\end{frame}

% =====================================================================
% 22 · CONCLUSIÓN Y LO QUE SIGUE
% =====================================================================
\begin{frame}
\tpsec{Conclusión}{}
\begin{lamina}
  \node[align=center,text=tpblue] at (8,1.5) {\lab{Conclusión}};
  \node[align=center,text=tpink,text width=14cm] at (8,2.6) {\bfseries\tpHone Le ganan a la mediana,\\[2pt]\hl{pero por poco}};
  \cardc{tpok}{0.9}{3.9}{4.55}{3.0}{{\bfseries\tpHthree\color{tpok}Hay señal}\par\vspace{4pt}\tpSmall Los tres modelos bajan el error: de 32,1 a 27,0 minutos de MAE (16 \%). El peso está en el tipo de hoja y en el primer ping.}
  \cardc{tpwarn}{5.75}{3.9}{4.55}{3.0}{{\bfseries\tpHthree\color{tpwarn}Falta mejorar}\par\vspace{4pt}\tpSmall El error sigue siendo grande, y en el 42 \% de los eventos no hay GPS previo.}
  \card{10.6}{3.9}{4.5}{3.0}{{\bfseries\tpHthree Lo que sigue}\par\vspace{4pt}\tpSmall Ajustar hiperparámetros, tratar los eventos de 4 h como censurados, separar desvíos cortos y largos, y confirmar con Tecfield lo de las hojas vencidas.}
  \path[fill=tpnavy,rounded corners=8pt] (0.9,7.2) rectangle (15.1,8.3);
  \node[align=center,text=white!80,font=\ttfamily\tpTiny] at (8,7.5) {\textls[60]{CÓDIGO, NOTEBOOK E INFORMES}};
  \node[align=center,text=white,font=\bfseries\tpHthree] at (8,7.9) {github.com/sebareibold/prediccion-duracion-desvios};
\end{lamina}
\note{\g{cierre}{conclusion} \g{cierre}{sigue}}
\end{frame}

\end{document}
